\chapter{Complexity — NCAC}
%%% generated by the blueprint pipeline from TCSlib.Complexity.CircuitComplexity.NCAC
\section{Overview}

Declarations of the Lean module \texttt{TCSlib.Complexity.CircuitComplexity.NCAC}.

\section{Declarations}

\begin{definition}[Languages in bounded-depth circuit families]
\lean{Language.InNC}
\label{Language.InNC}
\leanok
\uses{BoolCircuit.DAGCircuitFamily, BoolCircuit.DAGCircuitFamily.HasFaninTwo, BoolCircuit.DAGCircuitFamily.HasPolylogDepth, BoolCircuit.DAGCircuitFamily.IsPolySize, BoolCircuit.DAGCircuitFamily.language}
For a natural number $d$ and a language $L$ over Boolean words, say that $L$ belongs to
this class if there is a family of Boolean circuits represented as directed acyclic
graphs, with gates of fan-in at most two, polynomial size, and polylogarithmic depth
with parameter $d$, whose language is exactly $L$.
\end{definition}

\begin{definition}[Languages in a fixed AC level]
\lean{Language.InAC}
\statementsource{lmn-constant-depth-fourier}{pdf-63759bd909d5-p004-b007}
\label{Language.InAC}
\leanok
\uses{BoolCircuit.DAGCircuitFamily, BoolCircuit.DAGCircuitFamily.HasPolylogDepth, BoolCircuit.DAGCircuitFamily.IsPolySize, BoolCircuit.DAGCircuitFamily.IsWellFormed, BoolCircuit.DAGCircuitFamily.language}
For a natural number $d$, a language $L$ over the Boolean alphabet belongs to the $d$th
AC level if there is a well-formed family of directed acyclic circuit graphs computing
$L$, whose size is bounded by a polynomial in the input length and whose depth is
bounded by a polynomial in the logarithm of the input length, with the depth bound
determined by $d$.
\end{definition}

\begin{definition}[Circuit depth level]
\lean{BoolCircuit.NCLevel}
\label{BoolCircuit.NCLevel}
\leanok
\uses{Language.InNC}
For a natural number $d$, the depth-$d$ level is the set of languages over Boolean words
that are recognized by a family of directed acyclic Boolean circuits with fan-in at most
two, polynomial size, and polylogarithmic depth with parameter $d$.
\end{definition}

\begin{definition}[A fixed AC level]
\lean{BoolCircuit.ACLevel}
\statementsource{lmn-constant-depth-fourier}{pdf-63759bd909d5-p004-b007}
\label{BoolCircuit.ACLevel}
\leanok
\uses{Language.InAC}
For a natural number $d$, the $d$th AC level is the set of all languages over the
Boolean alphabet that belong to that level.
\end{definition}

\begin{definition}[The class NC]
\lean{BoolCircuit.NC}
\label{BoolCircuit.NC}
\leanok
\uses{BoolCircuit.NCLevel}
The class $\mathrm{NC}$ is the union, over all natural numbers $i\ge 1$, of the
depth-$i$ levels of languages recognized by polynomial-size, fan-in-two directed acyclic
circuit families with polylogarithmic depth.
\end{definition}

\begin{definition}[The AC class]
\lean{BoolCircuit.AC}
\statementsource{lmn-constant-depth-fourier}{pdf-63759bd909d5-p004-b007}
\label{BoolCircuit.AC}
\leanok
\uses{BoolCircuit.ACLevel}
The class of Boolean languages in AC is the union of the AC levels over all natural
numbers.
\end{definition}

\begin{theorem}[Characterization of NC]
\lean{BoolCircuit.mem_NC_iff}
\label{BoolCircuit.mem_NC_iff}
\leanok
\uses{BoolCircuit.NC, BoolCircuit.NCLevel, Language.InNC}
\difficulty{1}
For every language $L$ over the Boolean alphabet, $L$ belongs to $\mathrm{NC}$ if and
only if there is an integer $i\ge 1$ and a circuit family recognizing exactly $L$ such
that every circuit has fan-in at most two, the family has polynomial size, and its depth
is bounded by a polynomial in $\log n+1$ of degree $i$.
\end{theorem}

\begin{theorem}[Characterization of AC]
\lean{BoolCircuit.mem_AC_iff}
\label{BoolCircuit.mem_AC_iff}
\leanok
\uses{BoolCircuit.AC, BoolCircuit.ACLevel, Language.InAC}
\difficulty{1}
For every language $L$ over the Boolean alphabet, $L$ belongs to $\mathrm{AC}$ if and
only if there is a natural number $i$ such that $L$ has a well-formed, polynomial-size
circuit family whose depth is bounded by a polynomial in $\log n+1$ of degree $i$; the
family recognizes exactly $L$.
\end{theorem}

\begin{theorem}[NC circuits imply AC circuits]
\lean{Language.InNC.inAC}
\label{Language.InNC.inAC}
\leanok
\uses{BoolCircuit.DAGCircuitFamily.HasFaninTwo.isWellFormed, Language.InAC, Language.InNC}
\difficulty{4}
Let $d$ be a natural number and $L$ a language of Boolean strings. If $L$ is recognized
by a well-formed circuit family with fan-in at most two, polynomial size, and depth at
most $b(\lfloor\log_2 n\rfloor+1)^d$ on inputs of length $n$ for some natural number
$b$, then $L$ is recognized by a well-formed circuit family of polynomial size with the
same depth bound for some such $b$.
\end{theorem}

\begin{theorem}[Inclusion of Circuit Complexity Levels]
\lean{BoolCircuit.NCLevel_subset_ACLevel}
\label{BoolCircuit.NCLevel_subset_ACLevel}
\leanok
\uses{BoolCircuit.ACLevel, BoolCircuit.NCLevel, Language.InNC.inAC}
\difficulty{1}
For every natural number $i$, every Boolean language in the depth-$i$ class $	extsf{NC}$
also belongs to the depth-$i$ class $	extsf{AC}$.
\end{theorem}

\begin{theorem}[Polynomial input bound]
\lean{BoolCircuit.poly_add_input}
\label{BoolCircuit.poly_add_input}
\leanok
\difficulty{4}
For natural numbers $a,k,n,s$, if $s \le a(n+1)^k$, then $n+s \le (a+1)(n+1)^{k+1}$.
\end{theorem}

\begin{theorem}[Polylogarithmic bound after adding one]
\lean{BoolCircuit.polylog_add_one}
\label{BoolCircuit.polylog_add_one}
\leanok
\difficulty{3}
For natural numbers $b,d,L,t$, if $t \le b(L+1)^d$, then $t+1 \le (b+1)(L+1)^d$.
\end{theorem}

\begin{definition}[Tree circuit family as a DAG family]
\lean{BoolCircuit.TreeCircuitFamily.toDAG}
\label{BoolCircuit.TreeCircuitFamily.toDAG}
\leanok
\uses{BoolCircuit.DAGCircuitFamily, BoolCircuit.TreeCircuit.toDAG, BoolCircuit.TreeCircuitFamily}
Every family of Boolean circuit trees determines a family of directed acyclic Boolean
circuits by converting each circuit in the family, for every input length, to a directed
acyclic circuit.
\end{definition}

\begin{theorem}[Language preservation under tree-to-DAG conversion]
\lean{BoolCircuit.TreeCircuitFamily.language_toDAG}
\label{BoolCircuit.TreeCircuitFamily.language_toDAG}
\leanok
\uses{BoolCircuit.DAGCircuitFamily.language, BoolCircuit.DAGCircuitFamily.mem_language_iff, BoolCircuit.TreeCircuit.toDAG_eval, BoolCircuit.TreeCircuitFamily, BoolCircuit.TreeCircuitFamily.language, BoolCircuit.TreeCircuitFamily.mem_language_iff, BoolCircuit.TreeCircuitFamily.toDAG}
\difficulty{3}
For every family of Boolean tree circuits, the language recognized by its associated
directed acyclic graph circuit family is equal to the language recognized by the tree
circuit family.
\end{theorem}

\begin{theorem}[Polynomial size under DAG conversion]
\lean{BoolCircuit.TreeCircuitFamily.isPolySize_toDAG}
\label{BoolCircuit.TreeCircuitFamily.isPolySize_toDAG}
\leanok
\uses{BoolCircuit.DAGCircuitFamily.IsPolySize, BoolCircuit.TreeCircuit.toDAG_size_le, BoolCircuit.TreeCircuitFamily, BoolCircuit.TreeCircuitFamily.IsPolySize, BoolCircuit.TreeCircuitFamily.toDAG, BoolCircuit.poly_add_input}
\difficulty{4}
Let $C$ be a family of Boolean circuits indexed by input length. If there are natural
numbers $a$ and $k$ such that, for every $n$, the size of $C_n$ is at most $a(n+1)^k$,
then the family of directed acyclic graph circuits obtained from $C$ has polynomial
size: there are natural numbers $a'$ and $k'$ such that every circuit at input length
$n$ has size at most $a'(n+1)^{k'}$.
\end{theorem}

\begin{theorem}[Polylogarithmic depth under DAG conversion]
\lean{BoolCircuit.TreeCircuitFamily.hasPolylogDepth_toDAG}
\label{BoolCircuit.TreeCircuitFamily.hasPolylogDepth_toDAG}
\leanok
\uses{BoolCircuit.DAGCircuitFamily.HasPolylogDepth, BoolCircuit.TreeCircuit.toDAG_depth_le, BoolCircuit.TreeCircuitFamily, BoolCircuit.TreeCircuitFamily.HasPolylogDepth, BoolCircuit.TreeCircuitFamily.toDAG, BoolCircuit.polylog_add_one}
\difficulty{4}
Let $C$ be a family of Boolean circuits indexed by input length, and let $d$ be a
natural number. If there is a natural number $b$ such that, for every $n$, the depth of
the circuit $C_n$ is at most $b(\log_2 n+1)^d$, then the family of directed acyclic
graph circuits obtained from $C$ satisfies the same depth bound with some natural-number
constant.
\end{theorem}

\begin{theorem}[Tree classes lie in circuit classes]
\lean{Language.InTreeNC.inNC}
\label{Language.InTreeNC.inNC}
\leanok
\uses{BoolCircuit.TreeCircuit.toDAG_isFaninTwo, BoolCircuit.TreeCircuitFamily.hasPolylogDepth_toDAG, BoolCircuit.TreeCircuitFamily.isPolySize_toDAG, BoolCircuit.TreeCircuitFamily.language_toDAG, BoolCircuit.TreeCircuitFamily.toDAG, Language.InNC, Language.InTreeNC}
\difficulty{4}
For a natural number $d$ and a language $L$ over Boolean strings, if $L$ is recognized
by a family of Boolean circuit trees with fan-in at most two, polynomial size, and depth
at most $b(\lfloor\log_2 n\rfloor+1)^d$ at every input length $n$, for some natural
number $b$, then $L$ is recognized by a family of well-formed directed acyclic Boolean
circuits with fan-in at most two, polynomial size, and depth at most $b'(\lfloor\log_2
n\rfloor+1)^d$ at every input length $n$, for some natural number $b'$.
\end{theorem}

\begin{theorem}[Tree-circuit classes lie in circuit classes]
\lean{Language.InTreeAC.inAC}
\label{Language.InTreeAC.inAC}
\leanok
\uses{BoolCircuit.TreeCircuit.toDAG_isWellFormed, BoolCircuit.TreeCircuitFamily.hasPolylogDepth_toDAG, BoolCircuit.TreeCircuitFamily.isPolySize_toDAG, BoolCircuit.TreeCircuitFamily.language_toDAG, BoolCircuit.TreeCircuitFamily.toDAG, Language.InAC, Language.InTreeAC}
\difficulty{4}
For a natural number $d$ and a language $L$ over Boolean strings, if $L$ is recognized
by a family of Boolean circuit trees of polynomial size and depth at most
$b(\lfloor\log_2 n\rfloor+1)^d$ at every input length $n$, for some natural number $b$,
then $L$ is recognized by a family of well-formed directed acyclic Boolean circuits of
polynomial size and depth at most $b'(\lfloor\log_2 n\rfloor+1)^d$ at every input length
$n$, for some natural number $b'$.
\end{theorem}

\begin{definition}[DAG circuit family as a tree family]
\lean{BoolCircuit.DAGCircuitFamily.toTree}
\label{BoolCircuit.DAGCircuitFamily.toTree}
\leanok
\uses{BoolCircuit.DAGCircuit.toTree, BoolCircuit.DAGCircuitFamily, BoolCircuit.TreeCircuitFamily}
Every family of directed acyclic Boolean circuits determines a family of Boolean circuit
trees by unfolding each circuit in the family, for every input length, into a tree
circuit.
\end{definition}

\begin{theorem}[Tree conversion preserves the accepted language]
\lean{BoolCircuit.DAGCircuitFamily.language_toTree}
\label{BoolCircuit.DAGCircuitFamily.language_toTree}
\leanok
\uses{BoolCircuit.DAGCircuit.toTree_eval, BoolCircuit.DAGCircuitFamily, BoolCircuit.DAGCircuitFamily.IsWellFormed, BoolCircuit.DAGCircuitFamily.language, BoolCircuit.DAGCircuitFamily.mem_language_iff, BoolCircuit.DAGCircuitFamily.toTree, BoolCircuit.TreeCircuitFamily.language, BoolCircuit.TreeCircuitFamily.mem_language_iff}
\difficulty{3}
Let $C$ be a family of Boolean directed acyclic circuits such that, in every circuit,
each gate has distinct arguments and every NOT gate has exactly one argument. The
language accepted by the tree family obtained by unfolding these circuits equals the
language accepted by $C$.
\end{theorem}

\begin{theorem}[A power of three bounded by a square]
\lean{BoolCircuit.three_pow_log_le}
\label{BoolCircuit.three_pow_log_le}
\leanok
\difficulty{4}
For every nonnegative integer $n$, $3^{\log_2 n} \le (n+1)^2$, where $\log_2 n$ denotes
the integer logarithm.
\end{theorem}

\begin{theorem}[NC circuits have tree representations]
\lean{Language.InNC.inTreeNC_one}
\label{Language.InNC.inTreeNC_one}
\leanok
\uses{BoolCircuit.DAGCircuit.depth, BoolCircuit.DAGCircuit.toTree_depth_le, BoolCircuit.DAGCircuit.toTree_maxFanin_le, BoolCircuit.DAGCircuit.toTree_size_le, BoolCircuit.DAGCircuitFamily.HasFaninTwo.isWellFormed, BoolCircuit.DAGCircuitFamily.language_toTree, BoolCircuit.DAGCircuitFamily.toTree, BoolCircuit.TreeCircuit.size, BoolCircuit.three_pow_log_le, Language.InNC, Language.InTreeNC}
\difficulty{4}
Let $L$ be a language of Boolean strings. If $L$ is recognized by a family of
well-formed Boolean circuits with fan-in at most two, polynomial size, and depth at most
$b(\lfloor\log_2 n\rfloor+1)$ on inputs of length $n$ for some natural number $b$, then
$L$ is recognized by a family of Boolean circuit trees with maximum fan-in at most two,
polynomial size, and depth at most $b'(\lfloor\log_2 n\rfloor+1)$ for some natural
number $b'$.
\end{theorem}

\begin{theorem}[Depth-zero AC languages have shallow tree circuits]
\lean{Language.InAC.inTreeAC_zero}
\label{Language.InAC.inTreeAC_zero}
\leanok
\uses{BoolCircuit.DAGCircuit.args_length_le_size, BoolCircuit.DAGCircuit.depth, BoolCircuit.DAGCircuit.size, BoolCircuit.DAGCircuit.toTree_depth_le, BoolCircuit.DAGCircuit.toTree_size_le, BoolCircuit.DAGCircuitFamily.language_toTree, BoolCircuit.DAGCircuitFamily.toTree, BoolCircuit.TreeCircuit.size, Language.InAC, Language.InTreeAC}
\difficulty{4}
If a language over Boolean strings is recognized by a well-formed family of Boolean
circuits of polynomial size and depth at most a constant times $(\lfloor\log_2
n\rfloor+1)^0$ on inputs of length $n$, then it is recognized by a family of Boolean
circuit trees of polynomial size and depth at most a constant times $(\lfloor\log_2
n\rfloor+1)^0$.
\end{theorem}

\begin{theorem}[Depth-one circuit characterizations]
\lean{Language.inNC_one_iff}
\label{Language.inNC_one_iff}
\leanok
\uses{Language.InNC, Language.InNC.inTreeNC_one, Language.InTreeNC, Language.InTreeNC.inNC}
\difficulty{1}
For every Boolean language $L$, recognition by a family of Boolean circuit trees with
fan-in at most two, polynomial size, and depth at most $b(loor{\text{?}})$—more
precisely, depth at most $b(\log_2 n+1)$ for all input lengths $n$, for some natural
number $b$—is equivalent to recognition by a family of directed acyclic Boolean circuits
with fan-in at most two, polynomial size, and depth at most $b(\log_2 n+1)$ for all $n$,
for some natural number $b$.
\end{theorem}

\begin{theorem}[Depth-zero circuit class equivalence]
\lean{Language.inAC_zero_iff}
\label{Language.inAC_zero_iff}
\leanok
\uses{Language.InAC, Language.InAC.inTreeAC_zero, Language.InTreeAC, Language.InTreeAC.inAC}
\difficulty{1}
For every language $L$ of Boolean strings, $L$ is recognized by a well-formed family of
directed acyclic Boolean circuits of polynomial size and depth bounded by a constant
times $(\log_2 n+1)^0$ exactly when it is recognized by a family of Boolean circuit
trees with polynomial size and depth bounded by a constant times $(\log_2 n+1)^0$.
\end{theorem}

\begin{theorem}[Depth-zero AC circuits lie in NC]
\lean{Language.InAC.inNC_one_of_zero}
\label{Language.InAC.inNC_one_of_zero}
\leanok
\uses{Language.InAC, Language.InAC.inTreeAC_zero, Language.InNC, Language.InTreeAC.inTreeNC_succ, Language.InTreeNC.inNC}
\difficulty{1}
If a language over Boolean strings is recognized by a well-formed family of Boolean
circuits of polynomial size and depth at most a constant times $(\lfloor\log_2
n\rfloor+1)^0$ on inputs of length $n$, then it is recognized by a family of Boolean
circuits with at most two inputs per gate, of polynomial size and depth at most a
constant times $(\lfloor\log_2 n\rfloor+1)^1$.
\end{theorem}

\begin{theorem}[AC depth gives bounded-fanin NC depth]
\lean{Language.InAC.inNC_succ}
\label{Language.InAC.inNC_succ}
\leanok
\uses{BoolCircuit.DAGCircuit.binarize, BoolCircuit.DAGCircuit.binarize_depth_le, BoolCircuit.DAGCircuit.binarize_eval, BoolCircuit.DAGCircuit.binarize_isFaninTwo, BoolCircuit.DAGCircuit.binarize_size_le, BoolCircuit.DAGCircuit.depth, BoolCircuit.DAGCircuit.size, BoolCircuit.DAGCircuitFamily.mem_language_iff, BoolCircuit.clog_poly_le, Language.InAC, Language.InNC}
\difficulty{5}
Let $d$ be a natural number and $L$ a language over Boolean strings. If $L$ is
recognized by a well-formed family of Boolean circuits of polynomial size and depth at
most a constant times $(\lfloor\log_2 n\rfloor+1)^d$ on inputs of length $n$, then it is
recognized by a family of Boolean circuits with at most two inputs per gate, of
polynomial size and depth at most a constant times $(\lfloor\log_2 n\rfloor+1)^{d+1}$.
\end{theorem}

\begin{theorem}[Alternating circuit class inclusion]
\lean{BoolCircuit.ACLevel_subset_NCLevel_succ}
\label{BoolCircuit.ACLevel_subset_NCLevel_succ}
\leanok
\uses{BoolCircuit.ACLevel, BoolCircuit.NCLevel, Language.InAC.inNC_succ}
\difficulty{1}
For every natural number $i$, every Boolean language computed by a polynomial-size,
well-formed circuit family of depth at most a constant multiple of $(\log_2 n+1)^i$ is
computed by a polynomial-size circuit family whose gates have fan-in at most two and
whose depth is at most a constant multiple of $(\log_2 n+1)^{i+1}$.
\end{theorem}

\begin{theorem}[Equality of NC and AC]
\lean{BoolCircuit.NC_eq_AC}
\label{BoolCircuit.NC_eq_AC}
\leanok
\uses{BoolCircuit.AC, BoolCircuit.NC, BoolCircuit.mem_AC_iff, BoolCircuit.mem_NC_iff, Language.InAC.inNC_succ, Language.InNC.inAC}
\difficulty{4}
The class $	extsf{NC}$, defined as the union of the depth-$i$ classes $	extsf{NC}^i$
over all natural numbers $i \ge 1$, equals the class $	extsf{AC}$, defined as the union
of the depth-$i$ classes $	extsf{AC}^i$ over all natural numbers $i$.
\end{theorem}

\begin{theorem}[NC circuits have polynomial size]
\lean{Language.InNC.inPPoly}
\label{Language.InNC.inPPoly}
\leanok
\uses{Language.InNC, Language.InPPoly, Language.inPPoly_iff}
\difficulty{1}
Let $d$ be a natural number and $L$ a language of Boolean strings. If $L$ is recognized
by a family of well-formed circuits with fan-in at most two, polynomial size, and depth
at most $b(\lfloor\log_2 n\rfloor+1)^d$ on inputs of length $n$ for some natural number
$b$, then there are natural numbers $a,k$ such that $L$ is recognized by a well-formed
circuit family with fan-in at most two and size at most $a(n+1)^k$ on inputs of length
$n$.
\end{theorem}

\begin{theorem}[AC languages have polynomial circuit size]
\lean{Language.InAC.inPPoly}
\label{Language.InAC.inPPoly}
\leanok
\uses{Language.InAC, Language.InAC.inNC_succ, Language.InNC.inPPoly, Language.InPPoly}
\difficulty{1}
If a language over Boolean strings is recognized by a well-formed family of Boolean
circuits of polynomial size and polylogarithmic depth, then there are natural numbers
$a$ and $k$ such that the language is recognized by a family of circuits with at most
two inputs per gate whose size on inputs of length $n$ is at most $a(n+1)^k$.
\end{theorem}

\begin{theorem}[NC Has Polynomial Size]
\lean{BoolCircuit.NC_subset_PPoly}
\label{BoolCircuit.NC_subset_PPoly}
\leanok
\uses{BoolCircuit.NC, BoolCircuit.PPoly, BoolCircuit.mem_NC_iff, Language.InNC.inPPoly}
\difficulty{3}
Every language in $	extsf{NC}$ has polynomial circuit size: $	extsf{NC}$ is contained in
the class of languages recognized by fan-in-two circuit families whose size is bounded
by $a(n+1)^k$ for some natural numbers $a$ and $k$.
\end{theorem}
